% filepath: /Users/antoniomone/Desktop/PhD/PHD Projects/CoMI-IRL/comi-irl/docs/graph_clustering_methodology.tex
\documentclass[11pt,a4paper]{article}

% Packages
\usepackage[utf8]{inputenc}
\usepackage[T1]{fontenc}
\usepackage{amsmath,amssymb,amsthm}
\usepackage{algorithm}
\usepackage{algpseudocode}
\usepackage{booktabs}
\usepackage{graphicx}
\usepackage{hyperref}
\usepackage{enumitem}
\usepackage{geometry}
\usepackage{xcolor}
\usepackage{tcolorbox}

\geometry{margin=1in}

% Custom commands
\newcommand{\vect}[1]{\mathbf{#1}}
\newcommand{\mat}[1]{\mathbf{#1}}
\newcommand{\set}[1]{\mathcal{#1}}
\newcommand{\norm}[1]{\left\|#1\right\|}
\newcommand{\abs}[1]{\left|#1\right|}
\newcommand{\E}{\mathbb{E}}
\newcommand{\R}{\mathbb{R}}
\newcommand{\MI}{\mathrm{I}}

\title{\textbf{Adaptive Graph-Based Trajectory Clustering:\\A Methodology for Behavioral Mode Discovery in CoMI-IRL}}
\author{}
\date{}

\begin{document}

\maketitle

\begin{abstract}
We present an adaptive graph-based clustering method for discovering behavioral modes in trajectory embeddings learned by a contrastive behavior encoder. The method automatically detects whether the embedding space contains isolated islands (well-separated clusters) or forms a continuous manifold (overlapping clusters), and applies the appropriate clustering strategy. For isolated data, we use stable-$k$ connected component detection; for overlapping data, we perform a joint sweep over graph density ($k$) and Leiden resolution ($\gamma$), selecting the optimal configuration by modularity. We analyze the role of Jacobian-based behavioral features and demonstrate that they are largely redundant when the encoder is trained with global-local InfoMax and segment contrastive objectives, as these losses already capture the relevant control dynamics.
\end{abstract}

\tableofcontents
\newpage

%==============================================================================
\section{Related Work}
%==============================================================================

\subsection{Core Algorithm: Leiden Community Detection}

\begin{itemize}[leftmargin=*]
    \item \textbf{Traag, V. A., Waltman, L., \& Van Eck, N. J.} (2019). \textit{From Louvain to Leiden: guaranteeing well-connected communities.} Scientific Reports, 9(1), 5233.
    \begin{itemize}
        \item Introduces the Leiden algorithm as an improvement over Louvain
        \item Guarantees well-connected communities (no disconnected subcommunities)
        \item Optimizes modularity with a resolution parameter for multi-scale community detection
    \end{itemize}
\end{itemize}

\subsection{Graph Construction and Spectral Methods}

\begin{itemize}[leftmargin=*]
    \item \textbf{Von Luxburg, U.} (2007). \textit{A tutorial on spectral clustering.} Statistics and Computing, 17(4), 395-416.
    \begin{itemize}
        \item Comprehensive tutorial on $k$-NN graph construction for clustering
        \item Discusses graph Laplacians and their spectral properties
        \item Foundation for understanding graph-based clustering methods
    \end{itemize}
    
    \item \textbf{Zelnik-Manor, L., \& Perona, P.} (2004). \textit{Self-tuning spectral clustering.} NeurIPS.
    \begin{itemize}
        \item Adaptive scaling for affinity matrices
        \item Local neighborhood-based similarity computation
    \end{itemize}
\end{itemize}

\subsection{Graph Clustering on High-Dimensional Embeddings}

\begin{itemize}[leftmargin=*]
    \item \textbf{Becht, E., McInnes, L., Healy, J., et al.} (2019). \textit{Dimensionality reduction for visualizing single-cell data using UMAP.} Nature Biotechnology, 37(1), 38-44.
    \begin{itemize}
        \item Uses $k$-NN graphs combined with Leiden clustering for single-cell analysis
        \item Demonstrates effectiveness of graph-based clustering on learned embeddings
        \item Directly analogous to our trajectory embedding clustering approach
    \end{itemize}
    
    \item \textbf{Wolf, F. A., Angerer, P., \& Theis, F. J.} (2018). \textit{SCANPY: large-scale single-cell gene expression data analysis.} Genome Biology, 19(1), 15.
    \begin{itemize}
        \item Implements Leiden clustering on $k$-NN graphs of high-dimensional embeddings
        \item Scalable pipeline for clustering millions of data points
    \end{itemize}
\end{itemize}

\subsection{Resolution Parameter and Modularity}

\begin{itemize}[leftmargin=*]
    \item \textbf{Reichardt, J., \& Bornholdt, S.} (2006). \textit{Statistical mechanics of community detection.} Physical Review E, 74(1), 016110.
    \begin{itemize}
        \item Introduces resolution parameter interpretation in modularity optimization
        \item Connects community detection to statistical physics
    \end{itemize}
    
    \item \textbf{Fortunato, S., \& Barthélemy, M.} (2007). \textit{Resolution limit in community detection.} PNAS, 104(1), 36-41.
    \begin{itemize}
        \item Discusses fundamental resolution limits in modularity-based methods
        \item Motivates the need for multi-resolution analysis
    \end{itemize}
\end{itemize}

\subsection{Contrastive Learning and Mutual Information}

\begin{itemize}[leftmargin=*]
    \item \textbf{Hjelm, R. D., Fedorov, A., Lavoie-Marchildon, S., et al.} (2019). \textit{Learning deep representations by mutual information estimation and maximization.} ICLR.
    \begin{itemize}
        \item Deep InfoMax (DIM) framework for representation learning
        \item Maximizes mutual information between global and local features
        \item Foundation for our global-local InfoMax loss
    \end{itemize}
    
    \item \textbf{Chen, T., Kornblith, S., Norouzi, M., \& Hinton, G.} (2020). \textit{A simple framework for contrastive learning of visual representations.} ICML.
    \begin{itemize}
        \item SimCLR framework for contrastive self-supervised learning
        \item Temperature-scaled InfoNCE loss
        \item Basis for our instance contrastive objective
    \end{itemize}
    
    \item \textbf{Bardes, A., Ponce, J., \& LeCun, Y.} (2022). \textit{VICReg: Variance-invariance-covariance regularization for self-supervised learning.} ICLR.
    \begin{itemize}
        \item Variance-Invariance-Covariance regularization
        \item Prevents representation collapse without negative samples
        \item Informs our embedding stability constraints
    \end{itemize}
\end{itemize}

\subsection{Trajectory Clustering}

\begin{itemize}[leftmargin=*]
    \item \textbf{Lee, J. G., Han, J., \& Whang, K. Y.} (2007). \textit{Trajectory clustering: a partition-and-group framework.} SIGMOD, 593-604.
    \begin{itemize}
        \item Framework for trajectory segmentation and clustering
        \item Partition-and-group approach for handling variable-length trajectories
    \end{itemize}
    
    \item \textbf{Yuan, G., Sun, P., Zhao, J., Li, D., \& Wang, C.} (2017). \textit{A review of moving object trajectory clustering algorithms.} Artificial Intelligence Review, 47(1), 123-144.
    \begin{itemize}
        \item Comprehensive survey of trajectory clustering methods
        \item Covers distance-based, density-based, and partition-based approaches
    \end{itemize}
\end{itemize}

\subsection{Jacobian-Based Behavioral Analysis}

\begin{itemize}[leftmargin=*]
    \item \textbf{Todorov, E., \& Jordan, M. I.} (2002). \textit{Optimal feedback control as a theory of motor coordination.} Nature Neuroscience, 5(11), 1226-1235.
    \begin{itemize}
        \item Theoretical foundation for understanding control sensitivity in motor behaviors
        \item Connects action-state relationships to behavioral strategies
    \end{itemize}
    
    \item \textbf{Nguyen-Tuong, D., \& Peters, J.} (2011). \textit{Model learning for robot control: a survey.} Cognitive Processing, 12(4), 319-340.
    \begin{itemize}
        \item Survey of dynamics model learning for robotics
        \item Discusses Jacobian estimation and its role in control
    \end{itemize}
\end{itemize}

%==============================================================================
\section{Problem Setting}
%==============================================================================

Given a set of $N$ expert trajectories $\{\tau_1, \ldots, \tau_N\}$ from $K$ unknown behavioral modes, each trajectory is encoded by a learned behavior encoder into a $D$-dimensional embedding:
\begin{equation}
    \vect{z}_i = f_\theta(\tau_i) \in \R^D
\end{equation}
where $f_\theta$ is a transformer-based encoder with parameters $\theta$.

\textbf{Goal:} Discover $K$ behavioral modes (clusters) \textit{without knowing $K$ a priori}, handling two distinct data regimes:

\begin{table}[h]
\centering
\begin{tabular}{@{}lll@{}}
\toprule
\textbf{Regime} & \textbf{Structure} & \textbf{Example Environments} \\
\midrule
Isolated & Well-separated islands, no overlap & Pusher-v4, Reacher-v4 \\
Overlapping & Single connected manifold & Walker2d-v4 \\
\bottomrule
\end{tabular}
\caption{Two distinct embedding space regimes requiring different clustering strategies.}
\end{table}

%==============================================================================
\section{Behavior Encoder Architecture}
%==============================================================================

Before discussing clustering, we describe the encoder that produces the embeddings, as its training objectives are crucial for understanding why certain post-hoc features may be redundant.

\subsection{Input Representation}

Each trajectory $\tau = \{(\vect{s}_t, \vect{a}_t)\}_{t=0}^{T-1}$ is processed as follows:

\textbf{Random Fourier Features (RFF):} States and actions are projected using RFF embeddings:
\begin{align}
    \phi_{\text{state}}(\vect{s}) &= \sqrt{\frac{2}{m_s}} \cos(\mat{W}_s \vect{s} + \vect{b}_s) \\
    \phi_{\text{action}}(\vect{a}) &= \sqrt{\frac{2}{m_a}} \cos(\mat{W}_a \vect{a} + \vect{b}_a)
\end{align}
where $\mat{W}_s \sim \mathcal{N}(0, \sigma_s^2 \mat{I})$, $\mat{W}_a \sim \mathcal{N}(0, \sigma_a^2 \mat{I})$, and $m_s, m_a$ are the RFF dimensions.

\begin{tcolorbox}[colback=blue!5,colframe=blue!40!black,title=RFF Hyperparameters]
\begin{itemize}
    \item $m_s = 64$ (state RFF dimension)
    \item $m_a = 32$ (action RFF dimension)
    \item $\sigma_s = \sigma_a = 0.001$ (Walker2d, Pusher) or $0.01$ (Reacher)
\end{itemize}
Small $\sigma$ values capture high-frequency patterns in state-action sequences.
\end{tcolorbox}

\textbf{Typed Token Embeddings:} Each timestep produces a token with type information:
\begin{equation}
    \vect{h}_t = [\phi_{\text{state}}(\vect{s}_t); \phi_{\text{action}}(\vect{a}_t)] + \vect{e}_{\text{type}}
\end{equation}

\textbf{CLS Token:} A learnable $[\text{CLS}]$ token is prepended and serves as the trajectory-level embedding after transformer processing.

\subsection{Training Objectives}

The encoder is trained with a composite loss:
\begin{equation}
    \mathcal{L} = \alpha \mathcal{L}_{\text{contrastive}} + \beta \mathcal{L}_{\text{InfoMax}} + \gamma \mathcal{L}_{\text{segment}}
\end{equation}

\subsubsection{Instance Contrastive Loss ($\alpha$ term)}

Standard contrastive loss treating augmentations of the same trajectory as positives:
\begin{equation}
    \mathcal{L}_{\text{contrastive}} = -\log \frac{\exp(\text{sim}(\vect{z}_i, \vect{z}_i^+) / \tau)}{\sum_{j} \exp(\text{sim}(\vect{z}_i, \vect{z}_j) / \tau)}
\end{equation}

\subsubsection{Global-Local InfoMax Loss ($\beta$ term)}

Maximizes mutual information between the global CLS token and local (state, action) tokens:
\begin{equation}
    \mathcal{L}_{\text{InfoMax}} = -\frac{1}{T} \sum_{t=1}^{T} \log \sigma\left( D_\psi(\vect{z}_{\text{CLS}}, \vect{h}_t) \right) - \log \sigma\left( -D_\psi(\vect{z}_{\text{CLS}}, \vect{h}_t^-) \right)
\end{equation}
where $D_\psi$ is a discriminator network and $\vect{h}_t^-$ are negative samples from other trajectories.

\begin{tcolorbox}[colback=green!5,colframe=green!40!black,title=Key Insight: InfoMax Captures Dynamics]
The InfoMax objective forces the CLS embedding to encode information about \textbf{how actions relate to states over time}. This is precisely what Jacobian-based features attempt to capture manually.
\end{tcolorbox}

\subsubsection{Segment Contrastive Loss ($\gamma$ term)}

Ensures the global embedding is consistent with random sub-trajectory segments:
\begin{equation}
    \mathcal{L}_{\text{segment}} = -\log \frac{\exp(\text{sim}(\vect{z}_{\text{full}}, \vect{z}_{\text{seg}}) / \tau)}{\sum_{j} \exp(\text{sim}(\vect{z}_{\text{full}}, \vect{z}_{\text{seg},j}) / \tau)}
\end{equation}

\begin{tcolorbox}[colback=yellow!5,colframe=yellow!40!black,title=Key Insight: Segment Loss Captures Temporal Patterns]
By requiring consistency between full trajectories and random segments, the encoder learns \textbf{temporal dynamics and control patterns} that are invariant to trajectory length and starting point.
\end{tcolorbox}

\subsection{Default Hyperparameters}

\begin{table}[h]
\centering
\begin{tabular}{@{}ll@{}}
\toprule
\textbf{Parameter} & \textbf{Value} \\
\midrule
Embedding dimension $D$ & 32 \\
Transformer heads & 4 \\
Transformer layers & 2 \\
Hidden dimension & 1024 \\
Dropout & 0.1 \\
$\alpha$ (contrastive weight) & 0.5 \\
$\beta$ (InfoMax weight) & 1.0 \\
$\gamma$ (segment weight) & 0.5 \\
\bottomrule
\end{tabular}
\caption{Encoder training hyperparameters.}
\end{table}

%==============================================================================
\section{Embedding Normalization}
%==============================================================================

All embeddings are $L_2$-normalized to project them onto the unit hypersphere $\set{S}^{D-1}$:
\begin{equation}
    \hat{\vect{z}}_i = \frac{\vect{z}_i}{\norm{\vect{z}_i}_2 + \epsilon}
\end{equation}
where $\epsilon = 10^{-8}$ prevents division by zero.

\subsection{Rationale for Normalization}

\begin{enumerate}
    \item \textbf{Cosine similarity becomes dot product:}
    \begin{equation}
        \hat{\vect{z}}_i^\top \hat{\vect{z}}_j = \cos(\theta_{ij})
    \end{equation}
    where $\theta_{ij}$ is the angle between the original vectors.
    
    \item \textbf{Scale invariance:} Removes magnitude effects, focusing purely on directional similarity.
    
    \item \textbf{Bounded distances:} Cosine distance is bounded:
    \begin{equation}
        d_{\cos}(i, j) = 1 - \hat{\vect{z}}_i^\top \hat{\vect{z}}_j \in [0, 2]
    \end{equation}
\end{enumerate}

%==============================================================================
\section{$k$-NN Graph Construction}
%==============================================================================

We construct a weighted undirected graph $\set{G} = (\set{V}, \set{E}, \mat{W})$ encoding the local neighborhood structure of the embedding space.

\subsection{Step 1: Compute $k$-Nearest Neighbors}

For each embedding $\hat{\vect{z}}_i$, find the $k$ most similar embeddings using cosine distance:
\begin{equation}
    d_{\cos}(i, j) = 1 - \hat{\vect{z}}_i^\top \hat{\vect{z}}_j
\end{equation}

Let $\set{N}_k(i) = \{j_1, \ldots, j_k\}$ denote the indices of the $k$ nearest neighbors of point $i$ (excluding $i$ itself).

\begin{tcolorbox}[colback=blue!5,colframe=blue!40!black,title=Intuition: Effect of $k$]
\begin{itemize}
    \item \textbf{Small $k$ (5--10):} Very local connections, may fragment cohesive clusters
    \item \textbf{Large $k$ (50--100):} Denser connections, may connect distinct clusters
\end{itemize}
\end{tcolorbox}

\subsection{Step 2: Compute Edge Weights via RBF Kernel}

For each neighbor $j \in \set{N}_k(i)$, assign an edge weight using a Radial Basis Function (RBF) kernel on cosine similarity:
\begin{equation}
    w_{ij}^{\text{directed}} = \exp\left(\frac{\text{sim}(i,j)}{\sigma}\right) = \exp\left(\frac{\hat{\vect{z}}_i^\top \hat{\vect{z}}_j}{\sigma}\right)
\end{equation}
where $\sigma = 1.0$ is the kernel bandwidth.

\begin{tcolorbox}[colback=green!5,colframe=green!40!black,title=Intuition: RBF Weighting]
\begin{itemize}
    \item Higher similarity $\rightarrow$ Higher edge weight $\rightarrow$ Stronger connection
    \item Exponential transformation amplifies differences
    \item Creates a ``soft'' graph where edge strength reflects neighborhood confidence
\end{itemize}
\end{tcolorbox}

\subsection{Step 3: Symmetrization}

The $k$-NN relation is \textit{not symmetric}: if $j \in \set{N}_k(i)$, it does not guarantee $i \in \set{N}_k(j)$. We symmetrize by taking the element-wise maximum:
\begin{equation}
    w_{ij} = \max(w_{ij}^{\text{directed}}, w_{ji}^{\text{directed}})
\end{equation}

This ensures:
\begin{itemize}
    \item No information is lost from asymmetric neighborhoods
    \item The graph represents ``mutual similarity'' in a generous sense
    \item Isolated points are less likely to become disconnected
\end{itemize}

%==============================================================================
\section{Jacobian-Based Behavioral Features}
%==============================================================================

While trajectory embeddings capture high-level behavioral patterns, one might consider enhancing edge weights using features derived from the \textbf{action-state Jacobian}. We describe these features here but argue in Section~\ref{sec:redundancy} that they are largely redundant given our encoder's training objectives.

\subsection{Motivation}

Two trajectories may have similar embeddings but exhibit fundamentally different \textit{control strategies}:
\begin{itemize}
    \item One policy may use high-gain, reactive control
    \item Another may use smooth, anticipatory control
    \item They might reach similar states via different dynamical paths
\end{itemize}

The Jacobian $\frac{\partial \vect{s}_{t+1}}{\partial \vect{a}_t}$ captures \textit{how sensitive the next state is to the current action}---a signature of the underlying control strategy.

\subsection{Jacobian Estimation via Finite Differences}

For a trajectory $\tau = \{(\vect{s}_t, \vect{a}_t)\}_{t=0}^{T-1}$, the true Jacobian at time $t$ is:
\begin{equation}
    \mat{J}_t = \frac{\partial \vect{s}_{t+1}}{\partial \vect{a}_t} \in \R^{S \times A}
\end{equation}
where $S$ is the state dimension and $A$ is the action dimension.

Since we lack analytical access to environment dynamics, we estimate the Jacobian's magnitude using \textbf{finite differences}:

\textbf{State Change:}
\begin{equation}
    \Delta \vect{s}_t = \vect{s}_{t+1} - \vect{s}_t \in \R^S
\end{equation}

\textbf{Control Sensitivity (Jacobian Norm Proxy):}
\begin{equation}
    \gamma_t = \frac{\norm{\Delta \vect{s}_t}_2}{\norm{\vect{a}_t}_2 + \epsilon}
\end{equation}

This ratio measures how much state change results from a unit of action---a proxy for $\norm{\mat{J}_t}_F$ (the Frobenius norm of the Jacobian).

\subsection{Covariance-Based Jacobian Features}

For a more refined estimate, compute the \textbf{cross-covariance} between state changes and actions:
\begin{equation}
    \mat{C}_{\Delta s, a} = \text{Cov}(\Delta \vect{s}, \vect{a}) \in \R^{S \times A}
\end{equation}

Extract features via \textbf{Singular Value Decomposition (SVD)}:
\begin{equation}
    \mat{C}_{\Delta s, a} = \mat{U} \boldsymbol{\Sigma} \mat{V}^\top
\end{equation}
where $\boldsymbol{\Sigma} = \text{diag}(\sigma_1, \sigma_2, \ldots)$ with $\sigma_1 \geq \sigma_2 \geq \ldots$

\textbf{Key features:}
\begin{itemize}
    \item $\sigma_{\max} = \sigma_1$: Largest singular value (dominant control direction magnitude)
    \item $\rho = \frac{\sigma_1}{\sum_i \sigma_i}$: Singular value ratio (control dimensionality)
\end{itemize}

\subsection{Complete Behavioral Feature Vector}

For each trajectory $\tau_i$, we compute an 8-dimensional behavioral feature vector:
\begin{equation}
    \vect{b}_i = \begin{bmatrix}
        \bar{\gamma} \\
        \sigma_\gamma \\
        \gamma_{\max} \\
        \text{Var}_t(\gamma_t) \\
        \sigma_{\max} \\
        \rho \\
        \bar{a} \\
        \sigma_a
    \end{bmatrix}
    \quad
    \begin{array}{l}
        \text{(mean control sensitivity)} \\
        \text{(std of control sensitivity)} \\
        \text{(max control sensitivity)} \\
        \text{(temporal variance of sensitivity)} \\
        \text{(largest singular value of } \mat{C}_{\Delta s, a}) \\
        \text{(singular value concentration ratio)} \\
        \text{(mean action magnitude)} \\
        \text{(std of action magnitude)}
    \end{array}
\end{equation}

\begin{table}[h]
\centering
\begin{tabular}{@{}llp{7cm}@{}}
\toprule
\textbf{Feature} & \textbf{Symbol} & \textbf{Interpretation} \\
\midrule
Mean sensitivity & $\bar{\gamma}$ & Average control gain over trajectory \\
Sensitivity std & $\sigma_\gamma$ & Variability in control gain \\
Max sensitivity & $\gamma_{\max}$ & Peak control moment \\
Temporal variance & $\text{Var}_t(\gamma_t)$ & How much control gain changes over time \\
Max singular value & $\sigma_{\max}$ & Strength of dominant action-state coupling \\
SV ratio & $\rho$ & Dimensionality of control \\
Mean action & $\bar{a}$ & Average effort level \\
Action std & $\sigma_a$ & Effort variability \\
\bottomrule
\end{tabular}
\caption{Behavioral feature descriptions.}
\end{table}

\subsection{Edge Reweighting with Behavioral Similarity}

After computing behavioral features $\{\vect{b}_i\}_{i=1}^N$, normalize them:
\begin{equation}
    \tilde{\vect{b}}_i = \frac{\vect{b}_i - \bar{\vect{b}}}{\sigma_{\vect{b}} + \epsilon}
\end{equation}

For each edge $(i, j)$, compute \textbf{behavioral similarity} using an RBF kernel:
\begin{equation}
    s_{ij}^{\text{behav}} = \exp\left(-\frac{\norm{\tilde{\vect{b}}_i - \tilde{\vect{b}}_j}_2^2}{2\sigma_b^2}\right)
\end{equation}
where $\sigma_b = 1.0$ is the behavioral kernel bandwidth.

\textbf{Final edge weight} combines embedding-based and behavioral similarities:
\begin{equation}
    w_{ij}^{\text{final}} = (1 - \alpha) \cdot w_{ij}^{\text{emb}} + \alpha \cdot s_{ij}^{\text{behav}} \cdot w_{ij}^{\text{emb}}
\end{equation}
where $\alpha \in [0, 1]$ controls the influence of behavioral features (typically $\alpha = 0.2$ to $0.3$).

%==============================================================================
\section{On the Redundancy of Jacobian Features}
\label{sec:redundancy}
%==============================================================================

\begin{tcolorbox}[colback=red!5,colframe=red!40!black,title=Key Finding]
Given the encoder's training objectives (InfoMax + Segment Contrastive), Jacobian-based behavioral features are largely \textbf{redundant}. The encoder already learns to capture action-state dynamics through its loss functions.
\end{tcolorbox}

\subsection{What the Encoder Already Captures}

The three training objectives ensure that the CLS embedding encodes control dynamics:

\subsubsection{Global-Local InfoMax ($\beta$ term)}

\begin{equation}
    \mathcal{L}_{\text{InfoMax}} = -\MI(\vect{z}_{\text{CLS}}; \{\vect{h}_t\}_{t=1}^T)
\end{equation}

This loss \textbf{maximizes mutual information} between the global CLS token and local (state, action) tokens. For the CLS embedding to predict local tokens, it must encode:
\begin{itemize}
    \item How actions relate to states at each timestep
    \item The temporal structure of state-action sequences
    \item Control patterns that distinguish different policies
\end{itemize}

\textbf{This is precisely what $\gamma_t = \norm{\Delta \vect{s}_t} / \norm{\vect{a}_t}$ attempts to capture.}

\subsubsection{Segment Contrastive Loss ($\gamma$ term)}

By enforcing consistency between full trajectory embeddings and sub-trajectory segments:
\begin{equation}
    \vect{z}_{\text{full}} \approx \vect{z}_{\text{segment}}
\end{equation}

the encoder learns representations that are:
\begin{itemize}
    \item Invariant to trajectory length
    \item Sensitive to local control dynamics (since segments must match)
    \item Robust to temporal variations
\end{itemize}

\textbf{This captures the temporal variance $\text{Var}_t(\gamma_t)$ and smoothness of control.}

\subsubsection{RFF Embeddings}

The Random Fourier Feature embeddings with small bandwidth ($\sigma = 0.001$):
\begin{equation}
    \phi(\vect{s}) = \sqrt{\frac{2}{m}} \cos(\mat{W} \vect{s} + \vect{b}), \quad \mat{W} \sim \mathcal{N}(0, \sigma^{-2} \mat{I})
\end{equation}

capture \textbf{high-frequency patterns} in state-action sequences, including rapid changes in control sensitivity.

\subsection{Mapping Jacobian Features to Encoder Components}

\begin{table}[h]
\centering
\begin{tabular}{@{}lll@{}}
\toprule
\textbf{Jacobian Feature} & \textbf{Encoder Captures Via} & \textbf{Mechanism} \\
\midrule
$\bar{\gamma}$ (mean sensitivity) & InfoMax & CLS must encode action-state relationships \\
$\sigma_\gamma$ (sensitivity variance) & Segment contrastive & Temporal patterns in sub-trajectories \\
$\gamma_{\max}$ (max sensitivity) & Attention mechanism & Peak control moments get high attention \\
$\text{Var}_t(\gamma_t)$ & Segment contrastive & Consistency across trajectory segments \\
$\sigma_{\max}$ (SVD feature) & Cross-attention & State-action coupling in transformer \\
$\bar{a}, \sigma_a$ (action stats) & Action token embeddings & Direct encoding of action magnitudes \\
\bottomrule
\end{tabular}
\caption{How encoder training captures information equivalent to Jacobian features.}
\end{table}

\subsection{Empirical Verification}

To verify redundancy, compute the correlation between embedding-based and Jacobian-based similarity structures:

\begin{equation}
    \mat{S}^{\text{emb}}_{ij} = \hat{\vect{z}}_i^\top \hat{\vect{z}}_j, \quad
    \mat{S}^{\text{behav}}_{ij} = \exp\left(-\frac{\norm{\tilde{\vect{b}}_i - \tilde{\vect{b}}_j}^2}{2}\right)
\end{equation}

\begin{equation}
    r = \text{Corr}\left( \text{vec}(\mat{S}^{\text{emb}}), \text{vec}(\mat{S}^{\text{behav}}) \right)
\end{equation}

\textbf{Interpretation:}
\begin{itemize}
    \item $r > 0.7$: High redundancy — embeddings already capture behavioral dynamics
    \item $r < 0.4$: Low redundancy — Jacobian features add novel information
\end{itemize}

\subsection{Recommendation}

Given the sophisticated encoder training (InfoMax + Segment Contrastive + RFF), we \textbf{disable behavioral features by default}:

\begin{verbatim}
use_behavioral_features = False  # Default
\end{verbatim}

Retain the capability for:
\begin{itemize}
    \item \textbf{Ablation studies:} Quantify how much Jacobian features help/hurt
    \item \textbf{Cold-start scenarios:} When encoder training fails or data is scarce
    \item \textbf{Interpretability:} Jacobian features are human-readable
    \item \textbf{Cross-domain transfer:} Physics-based features may generalize better
\end{itemize}

%==============================================================================
\section{Automatic Regime Detection}
%==============================================================================

This is the \textbf{critical decision point} determining which clustering strategy to use.

\subsection{Connected Components}

A \textbf{connected component} of a graph is a maximal set of vertices with paths between every pair:
\begin{equation}
    \set{C} \subseteq \set{V} \text{ is a connected component if } \forall u, v \in \set{C}, \exists \text{ path } u \leadsto v \text{ in } \set{G}
\end{equation}

If a graph has multiple connected components, there are ``islands'' of vertices with \textit{no edges between them}.

\subsection{Detection Algorithm}

We check connectivity at five probe values: $k \in \{15, 30, 50, 75, 100\}$:

\begin{algorithm}[H]
\caption{Regime Detection}
\begin{algorithmic}[1]
\For{$k_{\text{probe}} \in \{15, 30, 50, 75, 100\}$}
    \State Build $k$-NN graph $\set{G}_{k_{\text{probe}}}$
    \State $n_{\text{comp}} \gets$ count connected components of $\set{G}_{k_{\text{probe}}}$
    \If{$n_{\text{comp}} \geq 2$}
        \State \Return \textsc{Isolated} regime
    \EndIf
\EndFor
\State \Return \textsc{Overlapping} regime
\end{algorithmic}
\end{algorithm}

\subsection{Rationale for Multiple Probe Values}

\begin{table}[h]
\centering
\begin{tabular}{@{}cl@{}}
\toprule
\textbf{Probe $k$} & \textbf{Purpose} \\
\midrule
15 & Catches well-separated islands \\
30 & Catches moderately separated islands \\
50 & Catches marginally separated islands \\
75, 100 & Final check for edge cases \\
\bottomrule
\end{tabular}
\end{table}

\textbf{Key insight:} If clusters are truly isolated, then \textit{for any reasonable $k$}, each cluster forms its own connected component. Increasing $k$ adds edges \textit{within} each cluster but never creates edges \textit{between} clusters.

%==============================================================================
\section{Isolated Regime: Stable-$k$ Connected Components}
%==============================================================================

When the embedding space contains isolated islands, we use connected components directly as clusters.

\subsection{The Stable-$k$ Algorithm}

\textbf{Goal:} Find the smallest $k$ where the connected component partition \textit{stabilizes}.

\begin{algorithm}[H]
\caption{Find Stable $k$ (with Early Exit)}
\begin{algorithmic}[1]
\Require Embeddings, $k_{\text{range}} = (k_{\min}, k_{\max})$, stability\_window $= 3$
\Ensure Optimal $k^*$, cluster labels
\State Sample $k$ values: $k_1, k_2, \ldots, k_n$ uniformly in $[k_{\min}, k_{\max}]$
\State stable\_count $\gets 0$
\For{each $k_i$}
    \State Build $k$-NN graph $\set{G}_{k_i}$
    \State Compute connected components: $\set{P}_{k_i} = \{C_1, C_2, \ldots, C_m\}$
    \If{$i > 1$}
        \State Compute $\text{ARI}(\set{P}_{k_i}, \set{P}_{k_{i-1}})$
        \If{$\text{ARI} = 1.0$}
            \State stable\_count $\gets$ stable\_count $+ 1$
            \If{stable\_count $\geq$ stability\_window}
                \State $k^* \gets k_{i - \text{stability\_window} + 1}$ \Comment{First $k$ in stable window}
                \State \Return $k^*$, $\set{P}_{k^*}$ \Comment{\textbf{Early exit}}
            \EndIf
        \Else
            \State stable\_count $\gets 0$
        \EndIf
    \EndIf
\EndFor
\State \Return heuristic selection if no stable window found
\end{algorithmic}
\end{algorithm}

\begin{tcolorbox}[colback=green!5,colframe=green!40!black,title=Efficiency: Early Exit]
Once stability is detected (3 consecutive $k$ values with identical partitions), the algorithm immediately returns. This avoids checking the remaining $k$ values, reducing computation from $O(n)$ to $O(\text{stable position})$.
\end{tcolorbox}

\subsection{Adjusted Rand Index (ARI)}

Connected components assign arbitrary integer labels. The \textbf{Adjusted Rand Index} compares partitions structurally:
\begin{equation}
    \text{ARI}(\set{P}_1, \set{P}_2) = \frac{\text{RI} - \E[\text{RI}]}{\max(\text{RI}) - \E[\text{RI}]}
\end{equation}
where RI counts pairs of points that are either together in both partitions or separate in both.

\textbf{ARI = 1.0} means the partitions are \textit{identical} (up to label permutation).

\subsection{Stability Window}

We require \textbf{3 consecutive stable $k$ values} before declaring stability, preventing:
\begin{itemize}
    \item False positives from coincidental matches
    \item Sensitivity to noise in the $k$-NN graph
\end{itemize}

%==============================================================================
\section{Overlapping Regime: Joint $k$-Resolution Sweep}
%==============================================================================

When the embedding space is a single connected manifold, we use the \textbf{Leiden algorithm} with joint parameter search.

\subsection{The Leiden Algorithm}

Leiden optimizes the \textbf{modularity function} with resolution parameter $\gamma$:
\begin{equation}
    Q_\gamma = \frac{1}{2m} \sum_{i,j} \left[ w_{ij} - \gamma \frac{s_i s_j}{2m} \right] \delta(c_i, c_j)
\end{equation}

\textbf{Components:}
\begin{itemize}
    \item $m = \frac{1}{2}\sum_{i,j} w_{ij}$: Total edge weight
    \item $s_i = \sum_j w_{ij}$: \textbf{Strength} (weighted degree) of node $i$
    \item $c_i \in \{1, \ldots, K\}$: Cluster assignment of node $i$
    \item $\delta(c_i, c_j)$: 1 if $c_i = c_j$, else 0
    \item $\gamma > 0$: \textbf{Resolution parameter}
\end{itemize}

\begin{tcolorbox}[colback=blue!5,colframe=blue!40!black,title=Intuition: Modularity]
\begin{itemize}
    \item $w_{ij}$: \textbf{Actual} edge weight between $i$ and $j$
    \item $\gamma \frac{s_i s_j}{2m}$: \textbf{Expected} edge weight under null model
    \item $Q_\gamma > 0$: Clusters have more internal edges than random $\rightarrow$ good partition
\end{itemize}
\end{tcolorbox}

\subsection{Effect of Resolution $\gamma$}

\begin{table}[h]
\centering
\begin{tabular}{@{}cll@{}}
\toprule
$\gamma$ & \textbf{Effect} & \textbf{Result} \\
\midrule
$< 1$ & Expected term smaller $\rightarrow$ easier to beat & Fewer, larger clusters \\
$= 1$ & Standard modularity & Balanced clusters \\
$> 1$ & Expected term larger $\rightarrow$ harder to beat & More, smaller clusters \\
\bottomrule
\end{tabular}
\end{table}

\subsection{Why Modularity for Selection (Not Silhouette)?}

\textbf{Observation:} In experiments, partitions with best NMI/ARI (matching ground truth) often had \textit{lower silhouette} than over-segmented partitions.

\textbf{Why this happens:}

Silhouette measures geometric compactness:
\begin{equation}
    S(i) = \frac{b(i) - a(i)}{\max(a(i), b(i))}
\end{equation}
where $a(i)$ = mean intra-cluster distance, $b(i)$ = mean distance to nearest other cluster.

\textbf{The problem:} Smaller clusters are more compact (lower $a(i)$) and further from others (higher $b(i)$), giving \textit{higher silhouette}---even if they represent over-segmentation.

\textbf{Why modularity works better:}
\begin{itemize}
    \item Modularity measures \textbf{graph-theoretic} community structure
    \item Over-segmentation $\rightarrow$ splits communities $\rightarrow$ reduces within-community edges $\rightarrow$ \textbf{lower modularity}
    \item Correct partition $\rightarrow$ maximizes edge ratio $\rightarrow$ \textbf{higher modularity}
\end{itemize}

\subsection{Joint Parameter Search}

We search over a 2D grid of $(k, \gamma)$ values:

\textbf{Default search grid:}
\begin{itemize}
    \item $k \in \{10, 15, 20, 25, 30, 40, 50, 70\}$
    \item $\gamma \in \{0.01, 0.025, 0.05, 0.1, 0.15, 0.2, 0.25, 0.3, 0.5, 0.7, 1.0, 1.3, 1.5, 2.0\}$
\end{itemize}

\begin{algorithm}[H]
\caption{Joint $k$-Resolution Sweep}
\begin{algorithmic}[1]
\For{each $k$ in $k\_values$}
    \State Build $k$-NN graph $\set{G}_k$
    \For{each $\gamma$ in $resolution\_values$}
        \State $\text{labels} \gets \text{Leiden}(\set{G}_k, \gamma)$
        \State Filter clusters with size $< n_{\min}$
        \If{valid\_clusters $\geq 2$}
            \State Store $(k, \gamma, \text{labels}, \text{modularity})$
        \EndIf
    \EndFor
\EndFor
\State $(k^*, \gamma^*) \gets \arg\max_{k,\gamma} \text{modularity}(k, \gamma)$
\State \Return labels from $(k^*, \gamma^*)$
\end{algorithmic}
\end{algorithm}

\subsection{Why Joint Search Over $k$ and $\gamma$?}

\begin{table}[h]
\centering
\begin{tabular}{@{}ll@{}}
\toprule
\textbf{Parameter} & \textbf{What it controls} \\
\midrule
$k$ & \textbf{Graph density}: How many neighbors define local similarity \\
$\gamma$ & \textbf{Cluster granularity}: How small clusters can be before merging \\
\bottomrule
\end{tabular}
\end{table}

These interact non-trivially:
\begin{itemize}
    \item Low $k$ + low $\gamma$: Sparse graph, few large clusters
    \item Low $k$ + high $\gamma$: Sparse graph, many tiny clusters (may fragment)
    \item High $k$ + low $\gamma$: Dense graph, clusters may merge inappropriately
    \item High $k$ + high $\gamma$: Dense graph, may find meaningful sub-communities
\end{itemize}

%==============================================================================
\section{Post-Processing}
%==============================================================================

After obtaining cluster assignments $\{c_1, \ldots, c_N\}$:

\subsection{Step 1: Filter Small Clusters}

Clusters with fewer than $n_{\min}$ points are relabeled as noise:
\begin{equation}
    c_i \leftarrow -1 \quad \text{if} \quad |\{j : c_j = c_i\}| < n_{\min}
\end{equation}

Default: $n_{\min} = \max(5, \lfloor 0.02 \cdot N \rfloor)$

\subsection{Step 2: Relabel Consecutively}

Valid cluster IDs are mapped to $\{0, 1, \ldots, K-1\}$.

\subsection{Step 3: Compute Cluster Centers}

For each cluster $c$, compute the normalized centroid:
\begin{equation}
    \boldsymbol{\mu}_c = \frac{\sum_{i: c_i = c} \hat{\vect{z}}_i}{\left\|\sum_{i: c_i = c} \hat{\vect{z}}_i\right\|_2}
\end{equation}

\subsection{Step 4: Build Cluster Registry}

For each cluster, store metadata:
\begin{itemize}
    \item $\boldsymbol{\mu}_c$: Normalized centroid
    \item $r_{95}$: 95th percentile cosine distance (cluster radius)
    \item count: Number of points
    \item $d_{\cos}^{\text{sorted}}$: Sorted cosine distances (for percentile queries)
\end{itemize}

%==============================================================================
\section{Complete Algorithm}
%==============================================================================

\begin{algorithm}[H]
\caption{Adaptive Graph-Based Trajectory Clustering}
\begin{algorithmic}[1]
\Require Embeddings $\mat{Z} \in \R^{N \times D}$, $k\_range$, $resolution\_values$, $n_{\min}$
\Ensure Cluster labels $\vect{c} \in \{-1, 0, \ldots, K-1\}^N$, cluster centers $\boldsymbol{\mu}$

\State \textbf{// Step 1: Normalize}
\State $\hat{\mat{Z}} \gets \mat{Z} / \norm{\mat{Z}}_2$

\State \textbf{// Step 2: Regime Detection}
\For{$k_{\text{probe}} \in \{15, 30, 50, 75, 100\}$}
    \State Build $k$-NN graph $\set{G}_{k_{\text{probe}}}$
    \State $n_{\text{comp}} \gets$ count connected components
    \If{$n_{\text{comp}} \geq 2$}
        \State regime $\gets$ \textsc{Isolated}
        \State \textbf{break}
    \EndIf
\EndFor
\If{all probes give $n_{\text{comp}} = 1$}
    \State regime $\gets$ \textsc{Overlapping}
\EndIf

\State \textbf{// Step 3: Clustering}
\If{regime = \textsc{Isolated}}
    \State Sweep $k$ to find stable partition via ARI
    \State Find $k^*$ where ARI $= 1.0$ for 3 consecutive $k$ (with early exit)
    \State \Return connected\_components($\set{G}_{k^*}$)
\Else
    \For{each $(k, \gamma)$ in grid}
        \State Build $k$-NN graph $\set{G}_k$
        \State labels $\gets$ Leiden($\set{G}_k$, $\gamma$)
        \State Store (labels, modularity)
    \EndFor
    \State Select $(k^*, \gamma^*)$ with highest \textbf{modularity}
    \State \Return Leiden partition at $(k^*, \gamma^*)$
\EndIf

\State \textbf{// Step 4: Post-process}
\State Filter clusters with size $< n_{\min} \rightarrow$ noise
\State Relabel to consecutive integers $\{0, 1, \ldots, K-1\}$
\State Compute normalized cluster centers
\end{algorithmic}
\end{algorithm}

\newpage
%==============================================================================
\section{Summary}
%==============================================================================

\begin{table}[h]
\centering
\begin{tabular}{@{}lll@{}}
\toprule
\textbf{Aspect} & \textbf{Isolated Regime} & \textbf{Overlapping Regime} \\
\midrule
Environments & Pusher-v4, Reacher-v4 & Walker2d-v4 \\
Embedding structure & Well-separated islands & Single connected manifold \\
Detection & $\geq 2$ connected components & 1 component at all probe $k$ \\
Clustering method & Connected components & Leiden algorithm \\
Parameter search & Sweep $k$ only (early exit) & Joint $(k, \gamma)$ sweep \\
Selection criterion & ARI stability & \textbf{Modularity} \\
\bottomrule
\end{tabular}
\caption{Comparison of clustering strategies for different data regimes.}
\end{table}

\subsection{On Jacobian Features}

\begin{table}[h]
\centering
\begin{tabular}{@{}lp{10cm}@{}}
\toprule
\textbf{Status} & \textbf{Recommendation} \\
\midrule
Default & \textbf{Disabled} --- encoder already captures control dynamics via InfoMax and segment contrastive losses \\
Enable when & Cold-start, few-shot learning, interpretability requirements, cross-domain transfer \\
\bottomrule
\end{tabular}
\end{table}

\subsection{Computational Complexity}

\begin{table}[h]
\centering
\begin{tabular}{@{}lll@{}}
\toprule
\textbf{Step} & \textbf{Complexity} & \textbf{Notes} \\
\midrule
Embedding normalization & $O(ND)$ & Linear in data size \\
$k$-NN search (brute) & $O(N^2 D)$ & Dominant for small $N$ \\
Connected components & $O(N + E)$ & Linear in graph size \\
Stable-$k$ (with early exit) & $O(n_{\text{stable}} \cdot N^2 D)$ & Often $n_{\text{stable}} \ll n_{\text{samples}}$ \\
Leiden clustering & $O(N \log N)$ & Per iteration \\
Joint $(k, \gamma)$ sweep & $O(n_k \cdot n_\gamma \cdot N \log N)$ & $\sim$112 Leiden runs \\
\bottomrule
\end{tabular}
\caption{Computational complexity of each algorithm step.}
\end{table}

%==============================================================================
\section{Implementation Details}
%==============================================================================

\subsection{Default Parameters}

\begin{table}[h]
\centering
\begin{tabular}{@{}lll@{}}
\toprule
\textbf{Parameter} & \textbf{Default} & \textbf{Description} \\
\midrule
\texttt{k\_range} & $(10, 100)$ & Range of $k$ values to search \\
\texttt{n\_samples} & 50 & Number of $k$ values to sample \\
\texttt{stability\_window} & 3 & Consecutive stable partitions required \\
\texttt{min\_cluster\_size} & $\max(5, 2\%N)$ & Minimum cluster size \\
\texttt{use\_behavioral\_features} & \texttt{False} & Jacobian feature reweighting \\
\texttt{behavioral\_alpha} & 0.3 & Weight for behavioral features (if enabled) \\
\texttt{seed} & 42 & Random seed for reproducibility \\
\bottomrule
\end{tabular}
\caption{Default clustering parameters.}
\end{table}

\subsection{Environment-Specific Settings}

\begin{table}[h]
\centering
\begin{tabular}{@{}llll@{}}
\toprule
\textbf{Environment} & \textbf{Expected Regime} & \textbf{$K$ (modes)} & \textbf{Notes} \\
\midrule
Pusher-v4 & Isolated & 6 & Well-separated goal positions \\
Reacher-v4 & Isolated & 6 & Well-separated target angles \\
Walker2d-v4 & Overlapping & 3 & Continuous gait variations \\
\bottomrule
\end{tabular}
\caption{Environment characteristics and expected clustering behavior.}
\end{table}

\end{document}